\documentclass[11pt]{article}
\usepackage[margin=25mm]{geometry}
\usepackage{fontspec}
\setmainfont{Latin Modern Roman}
\usepackage{amsmath,amssymb,mathtools}
\usepackage{unicode-math}
\setmathfont{Latin Modern Math}
\usepackage{xurl,hyperref}
\hypersetup{colorlinks=true,urlcolor=blue}
\setcounter{secnumdepth}{0}
\setlength{\parindent}{0pt}
\setlength{\parskip}{5pt}
\setlength{\emergencystretch}{3em}
\newcommand{\sref}[2]{\hyperlink{source:#1:#2}{[S#2]}}
\newcommand{\cataloguescope}[1]{\paragraph{Recorded target.}#1}
\newcommand{\R}{\mathbb R}
\newcommand{\C}{\mathbb C}
\newcommand{\N}{\mathbb N}
\newcommand{\Z}{\mathbb Z}
\newcommand{\Q}{\mathbb Q}
\newcommand{\F}{\mathbb F}
\newcommand{\eps}{\varepsilon}
\begin{document}
\section{Long-time behavior with non-generic initial data}\label{statement}

\subsection{Definitions and mathematical statement}

\paragraph{Shared notation.}$\mathbb N=\{1,2,\ldots\}$; the standard number fields have their usual meanings. All additional conventions are stated locally.


\paragraph{Definitions and notation.}For focusing NLS the associated AKNS scattering operator is non-self-adjoint. Discrete eigenvalues off the real axis correspond to soliton data. Accumulation of infinitely many such eigenvalues at the real continuous spectrum creates exceptional scattering singularities. For the Boussinesq equation the Lax operator is third order, and non-generic spectral behavior can likewise obstruct asymptotic analysis.

\paragraph{Catalogue wording retained for traceability.}
The long-time behavior of the solution of the Cauchy problem for a great many integrable systems on the line is now well-understood, using, for example, the Riemann-Hilbert/steepest descent method. The method depends on full knowledge of the nature of the spectrum of the associated Lax operator. For systems such as KdV, MKdV, defocusing NLS and the Toda lattice, for example, one is able to describe the solution asymptotically in complete detail for general initial data. But for focusing NLS, where the associated AKNS Lax operator $T$ is non-self-adjoint (here we make contact with Problem 7), the situation is different. For generic $T$ (i.e. an open dense set of $T$'s in any reasonable topology) the spectrum consists of the real line, where the spectrum is absolutely continuous, together with a finite number of simple eigenvalues (corresponding to solitons) off the real axis. The analysis of the long-time behavior of focusing NLS with such generic initial data proceeds in a straightforward manner similar to KdV, MKdV, etc. For general initial data, however, the situation is more complicated. For example, let $z_0>0$ be any positive number and let $D$ be any arbitrarily small open disk in the complex plane centered at $z_0$, such that $D\subset\{z\,:\,\mathbb{R}e z>0\}$. Let $D^+$ denote the intersection of $D$ with the upper half plane, and let $u(z)$ be an arbitrary function analytic in $D^+$, and continuous in $\overline{D^+}$. Let $B=\{z\in D^+\,:\, u(z)=0\}$. Then there exists (Zhou) an AKNS operator $T$ with infinitely smooth, rapidly decaying coefficients with the property that each point in $B$ is an $L^2(\mathbb{R})$-eigenvalue of $T$. In other words, there exist (non-generic) operators $T$ with Schwartz space coefficients which have $L^2$ spectrum accumulating on the real line at an essentially arbitrary rate. It is a very interesting question to determine what effect such singularities would have on the long-time behavior of the solution of NLS. In particular, recalling that focusing NLS provides a model for data transmission along communication cables, are the effects measurable? The difficulty that we encounter here is not limited to situations where the associated Lax operator is non-self-adjoint. Even in situations when the associated operator is self-adjoint, but of order greater than two, similar difficulties can arise. This is true, in particular, for the Boussinesq equation, where the associated Lax operator is third order. The long-time behavior of the solutions of the Boussinesq equation with general initial data is a very interesting problem with many challenges. Even in the case with generic initial data the situation is only partially understood.
\paragraph{Scope and source qualification.}The generic-data results and exceptional-data examples in the original text remain context. One must cover the specified resonant cases rather than substitute the regular reflection-coefficient assumptions of the established method.
\subsection{Short English statement}
Determine long-time behavior of integrable PDEs for exceptional initial data that produce singular scattering behavior.
\subsection{Research attempt}
\paragraph{Model, date and investigated scope.}
OpenAI GPT-6 Astra (configured model \texttt{gpt-6-astra}), ultra reasoning; generated 11 October 2026. Proved a scalar accumulation model has algebraic decay although each finite spectral truncation decays exponentially. Actual exceptional NLS and Boussinesq asymptotics remain unproved.

\paragraph{Sources and literature check.}
Read the supplied complete definition, catalogue wording, scope qualifications and source list for this canonical ID. The original-source citations are inherited as supplied references; no fresh primary-source or novelty search was performed for this attempt. No current open-status claim is inferred from this calculation.

\paragraph{Testing accumulation through exponential sums.}A basic obstacle in infinitely many near-continuous spectral contributions is failure of a uniform exponential time scale. For a transparent scalar analogue, take \(0<q<1\) and
\[F_q(t)=\sum_{j=1}^\infty q^j e^{-tq^j},\qquad t\ge0.
\]
This series converges absolutely and \(F_q(t)\to0\) by dominated convergence, because \(\sum q^j<\infty\). Nevertheless its decay is only of order \(1/t\). For \(t\ge q^{-1}\), choose \(m\ge1\) with \(q^{-m}\le t<q^{-m-1}\). The tail \(j\ge m+1\) has \(tq^j<1\), so
\[F_q(t)\ge e^{-1}\frac{q^{m+1}}{1-q}\ge\frac{e^{-1}q}{(1-q)t}.\]
For the upper bound, the same tail is at most \(q^{m+1}/(1-q)\le1/((1-q)t)\). For \(j=m-k\), \(k\ge0\), write \(q^j e^{-tq^j}=t^{-1}s_ke^{-s_k}\), where \(q^{-k}\le s_k<q^{-k-1}\). The sum is bounded by \(t^{-1}\sum_{k\ge0}q^{-k-1}e^{-q^{-k}}\), a finite constant times \(t^{-1}\). Thus \(c_q/t\le F_q(t)\le C_q/t\).
\paragraph{What fails in a finite-mode argument.}Every fixed finite partial sum decays exponentially; no fixed truncation gives the asymptotic order of the entire sum. The truncation error is small in absolute size but can dominate the retained terms at long times. This directly tests a tempting exchange of spectral truncation and time asymptotics.
\paragraph{Missing nonlinear bridge.}No identification of this scalar series with a focusing-NLS scattering reconstruction is asserted. The actual residue strengths, spectral singularities, nonlinear interactions and uniform reconstruction estimates must be controlled before it predicts a PDE tail or a measurable signal. Such control is the first missing step. The bundled third-order Boussinesq general-data problem is not addressed by this toy series either. This is an explicit countercheck of a proposed proof strategy, not a PDE long-time theorem.

\paragraph{Verification and adversarial audit.}
Checked the geometric-index choice, lower and upper constants, convergence and long-time/truncation noncommutation; no scalar-to-PDE identification was assumed. This was an LLM adversarial audit, not independent human verification; no theorem prover was used. Prompts A then B applied to this ID: Checked the geometric-index choice, lower and upper constants, convergence and long-time/truncation noncommutation; no scalar-to-PDE identification was assumed. Scope remains limited: Derive uniform nonlinear scattering reconstruction estimates for infinitely accumulating spectral data and justify any truncation limits.

Coordinator checked both geometric-tail bounds and the finite constant in the head estimate giving inverse-time decay. The nonlinear reconstruction bridge remains missing. This additional check was model-assisted, not independent human verification.

\paragraph{Tools, process and contributions.}
Python 3 and local shell used to read supplied files and serialize this manuscript; mathematics derived symbolically in the text, with no numerical experiment or formal proof assistant. Contribution prompt A was followed by prompt B. GPT-6 Astra with ultra reasoning drafted the research and performed its stated self-audit; any additional coordinator contribution is identified in the verification and audit. A GPT-6-based Codex coordinator assembled the manuscript and checked repository integration; its exact variant and reasoning configuration were not exposed. No human mathematical contribution, independent human proof check or formal proof certification is supplied. Complete derivations and the consulted source references are included here; private reasoning transcripts are not published.

\paragraph{Final status and first remaining gap.}
\textbf{UNRESOLVED}. Derive uniform nonlinear scattering reconstruction estimates for infinitely accumulating spectral data and justify any truncation limits.

\section{Completion Estimate}
\noindent\textbf{COMPLETION: 0\%}
The calculation is an elementary benchmark or conditional reduction; it establishes no claimed new portion of the complete source target.

\subsection{Sources}
\begin{itemize}
\item[\hypertarget{source:statement:1}{S1}] ULAM.AI and the UnsolvedMath Contributors, UnsolvedMath, catalogue entry. Catalogue curation under CC BY 4.0; definitions and English exposition added locally.
\url{https://huggingface.co/datasets/ulamai/UnsolvedMath}.
\item[\hypertarget{source:statement:2}{S2}] Percy Deift, Some open problems in random matrix theory and the theory of integrable systems (2007), arXiv:0712.0849, Problem 8.
\url{https://arxiv.org/abs/0712.0849}.
\end{itemize}
\label{end:statement}
\end{document}
